%%%
%
% $Autor: Wings $
% $Datum: 2021-05-14 $
% $Pfad: GitLab/MLEdgeComputer $
% $Dateiname:  JSTConnector
% $Version: 4620 $
%
% !TeX spellcheck = de_DE/GB
% !TeX program = pdflatex
% !BIB program = biber/bibtex
% !TeX encoding = utf8
%
%%%

\chapter{Steckverbinder vom Typ JST}

% https://en.wikipedia.org/wiki/JST_connector

\section{Steckverbinder-Serie}

JST-Steckverbinder sind elektrische Steckverbinder, die nach den ursprünglich von der J.S.T. Mfg. Co. (Japan Solderless Terminal) entwickelten Designstandards gefertigt werden \cite{JST:2016}. JST stellt zahlreiche Serien (Familien) und Rasterma\ss e (Abstand von Pin zu Pin) von Steckverbindern her.


JST fertigt eine große Anzahl von Steckverbinder-Serien (Familien). Die Leiterplattensteckverbinder (Draht-zu-Platine) sind in Ausführungen mit oberem (vertikalem) oder seitlichem (horizontalem) Einlass sowie als THT- (Through-Hole Technology) oder SMT- (Surface-Mount Technology) Varianten erhältlich \cite{JST:2020}. 

Diese Beschreibung bezieht sich ausschließlich auf die mechanische Struktur und das elektrische Verhalten. Die Pinbelegung ist nicht standardisiert.

Achtung: Der Begriff ``JST'' wird manchmal fälschlicherweise als umgangssprachlicher Ausdruck verwendet, der jeden kleinen wei\ss en elektrischen Steckverbinder bezeichnet, der auf Printed Circuit Board (PCB) montiert ist.




\section{JST VH}

%\URL{https://www.jst-mfg.com/product/index.php?series=262}

Die technischen Daten der JST VH-Serie sind in diesem Abschnitt angegeben. Sie stammen aus dem Datenblatt der Serie \cite{JST:2021}. 

\subsection{Product Profile}


\begin{tabular}{l|l}
	\centering
  Parameter & Profile \\ \hline
  Serie   & VH-Stecker \\ \hline
  Kategorie & Crimp-Steckverbinder (Wire-to-Board-Typ) \\ \hline
  Typ     & Crimpausführung, Kompakttyp, \\ 
           & Mit Sperrvorrichtung\\
           &  Abtrennbarer Typ \\ \hline
  Merkmal  & Mit sicherer Verriegelung\\ 
\end{tabular}
\captionof{table}{Product Profile}
\subsection{Spezifikation}


\begin{tabular}{l|l}
	\centering
  Parameter      & Spezifikation \\ \hline   	
  Stiftabstand        & 3.96mm \\ \hline
  Stiftreihen       & 1 \\ \hline
  Strombelastbarkeit & 10A (AWG\#16) \\ \hline
  Spannungsfestigkeit & 250V\\ \hline
  
 Leiterplattenmontagerichtung & Obereintritt, Seiteneintritt 
\end{tabular}  
\captionof{table}{Spezifikation}

\begin{center}
    \includegraphics[width=0.6\textwidth]{JSTConnectors/JSTVH}
    \captionof{figure}{JST connector type VH}
    \footnotesize \textbf{Reference:} \URL{https://www.jst-mfg.com/product/index.php?series=262}
    
\end{center}


\section{JST PH}

In diesem Abschnitt werden die technischen Daten der JST PH-Serie angegeben. Sie sind dem Datenblatt der Serie entnommen \cite{JST:2022}. 
%https://www.jst-mfg.com/product/index.php?series=199

\subsection{Product Profile}

\begin{center}
\begin{tabular}{l|l}
   Parameter & Profile \\ \hline
  Serie   & PH-Stecker \\ \hline
  Kategorie & Crimp-Steckverbinder (Wire-to-Board-Typ)\\ \hline
  Typ    & Crimp-Stil, Dünn-Typ \\
           & Trennbarer Typ \\
\end{tabular}
\captionof{table}{Product Profile}
\end{center}
\subsection{Spezifikation}

\begin{center}
\begin{tabular}{l|l}
	\centering
  Parameter      & Spezifikation \\ \hline	
  Stiftabstand  & 2.0mm  \\ \hline
  Strombelastbarkeit & 2A (AWG\#24) \\ \hline
  Spannungsfestigkeit & 100V \\ \hline
  Leiterplattenmontagerichtung & Obereintritt, Seiteneintritt \\ 
\end{tabular}
\captionof{table}{Spezifikation}
\end{center}
\begin{center}
    \includegraphics[width=0.6\textwidth]{JSTConnectors/JSTPH}
    \captionof{figure}{JST connector type VH } 
    \footnotesize \textbf{Reference:} \cite{JST:2022}
\end{center}

\subsection{JST PHR-2}

Eine Variante der Steckverbinderserie ist der Typ JST PHR-2, der 2 Stifte hat. Abbildung \ref{JST-2022} zeigt einen JST PHR-2 Steckverbinder.

\begin{center}
    \includegraphics[width=0.5\textwidth]{JSTConnectors/PHR-2.jpg}
    \captionof{figure}{JST-Stecker Typ PHR-2 }\label{JST-2022}
    \footnotesize \textbf{Reference:} \cite{JST:2022}
\end{center}

